\section{Durchführung}

\subsection{Slowkoinzidenzkreis}
Als erstes soll der Slowkreis aufgebaut und eingesellt werden. Dazu wird als erstes die Hochspannung des Detektors eingestellt. 
Zu beginn des Versuches waren diese schon eingeschaltet und eingestellt auf $\SI{610}{\kilo\electronvolt}$ und $\SI{108(1)}{\mikro\electronvolt}$
für den linken Detektor und $\SI{668}{\kilo\electronvolt}$ und $\SI{118(1)}{\mikro\electronvolt}$ für den rechten Detektor.
\subsubsection*{Niedervoltspannung}
Neben der Hochspannung muss auch die Niedervoltspannung vermessen werden.
Die gemessenen Werte sind in der Tabelle \ref{tab:niedervolt} aufgelistet, Außerdem sind die aufgenommen Oszillatorgraphen in dem Anhang unter Kapitel
\ref{sec:anhang} zu finden.

Die Niedervoltspannung wurde mit Hilfe einer speziellen Messspitze an den Oszilloskop angeschlossen und vermessen.
\begin{table}
\centering
\begin{tabular}{|c|c|c|}
\hline
    & Niedervoltkanel & Spannung [V] \\ \hline
    \multirow{7}{*}{links} & Grd & 0 \\ \hline
                        & +24 & 22 \\ \hline
                        & -24 & -22 \\ \hline
                        & +12 & 12 \\ \hline
                        & -12 & -12 \\ \hline
                        & +6 & unid. \\ \hline
                        & -6 & unid. \\ \hline
    \multirow{7}{*}{rechts} & Grd & 0.00 \\ \hline
                        & +24 & 23 \\ \hline
                        & -24 & -24 \\ \hline
                        & +12 & 12 \\ \hline
                        & -12 & -12 \\ \hline
                        & +6 & 6 \\ \hline
                        & -6 & -6 \\ \hline
\end{tabular}
\caption{Niedervoltspannungen der Photomultiplier.}
\label{tab:niedervolt}
\end{table}
Wie in der Tabelle zu sehen ist sind die $\SI{6}{\volt}$ Spannungen auf der linken Seite nicht lesbar, 
da hier mehrere Linien zu sehen sind von denen keiner zu den erwarteten 6 Volt passt.
Dies ist auch der Grund wieso mit dem Oszilloskop gemessen wurde, um genau solche Schwankungen und unstimmigkeiten
der Niedervoltspannungen zu erkennen.

\subsubsection*{Slow-pulse des Photomultipliers}
Um die Slow-pulse des Photomultipliers zu analysieren wurde zunächst das Signal nach dem Vorverstärker und das Signal
nach dem Hauptverstärker an ein Oszilloskop angeschlossen, der Aufbau der Schaltung ist in der Abbildung \ref{}
dargestellt.

\begin{figure}
    \centering
    \includegraphics[width=0.8\textwidth]{Koinzidenz}
    \caption{Aufbau der Schaltung zur Messung der Slow-pulse des Photomultipliers.}
    \label{fig:slowpulse}
\end{figure}
\subsubsection*{Triggerung mit dem SCA}
\subsubsection*{Energiespektrum der $^{22}$Na-Quelle}
\subsubsection*{Energiespektren der $^{22}$Na- und $^{133}$Ba-Quellen}
\subsubsection*{Einkanalfenster der $^{22}$Na-Quelle} 
\subsubsection*{Slowkoinzidenz herstellen}

\subsection{Fastkoinzidenzkreis}
\subsubsection*{Fast-pulse des Photomultipliers}
\subsubsection*{CFD-Diskriminatorschwelle}
\subsubsection*{Fastkoinzidenz herstellen}
\subsubsection*{Zeitlicher Abgleich der Fast- und Slowkoinzidenz}

\subsection{Zeitkalibration des TAC}
\subsubsection*{Aufnahme der Prompt-Kurve}
\subsubsection*{Bestimmung der Zeitauflösung und der Zeitkalibration}


\subsection{Lebensdauerbestimmung}
\subsubsection*{Einkanalfenser der $^{133}$Ba-Quelle}
\subsubsection*{Koinzidenzen kontrollieren}
\subsubsection*{Messung der Lebensdauerkurve}